\subsection{Interoperability — what!?}\label{subsec:Definitions}
\begin{frame}{Interoperability — what!?}
  \begin{definition}{Merriam Webster}
    \small \textit{in·ter·op·er·a·bil·i·ty} — ability of a system (such as a weapons system) to work with or use the parts or equipment of another system.
  \end{definition}

  \begin{definition}{Cambridge Dictionary}
    \small Interoperability: the degree to which two products, programs, etc.\ can be used together, or the quality of being able to be used together.
  \end{definition}

  \vspace{0.3cm}
  \begin{center}
    {\large As a data architect, I love me some well-defined terms!}
  \end{center}
\end{frame}

\begin{frame}{Interoperability — what!? (2)}
  \begin{definition}{U.S. Department of Energy}
    \small Interoperability: the ability of systems, personnel, or organizations to provide services to and accept services from other systems, personnel, or organizations, and to use the services so exchanged so that these organizations can operate together effectively.
  \end{definition}

  \vspace{0.3cm}
  \begin{definition}{Red Hat}
    \small Interoperability: the connection of data, applications, APIs, and devices across your IT organization to be more efficient, productive, and agile.
  \end{definition}
\end{frame}
